% Validation notes moved from the body during the presentation pass.
% Every number is produced by the same slots as the body.
\subsection*{Validation notes moved from the body}

\noindent\textbf{Part I.} The split is by environment, so every test number here is a cross-environment number, and
the threshold sweep reports retention, recall and precision on one curve, which makes the
precision collapse at high thresholds a stated property, not a surprise. Stressing the
two test environments separately (Appendix~\ref{app:ablation}): \texttt{\slotPEnvOneName}
holds \slotPEnvOneN{} frames with \slotPEnvOnePos{} positives and \texttt{\slotPEnvTwoName}
\slotPEnvTwoN{} frames with \slotPEnvTwoPos{}, scoring \slotPEnvOneAUC{} and
\slotPEnvTwoAUC{}. The pre-declared mark was AUC above \slotPassAUC{} with average precision
above prevalence: the pooled split passes, the courier-station environment does not.

\noindent\textbf{Part II.} Validation has three legs. The predictors are compared on identical windows, so
the comparison is paired by construction. Coverage is checked against the nominal
level, and the pass mark is coverage of at least \slotConfLevel{}\% minus a
tolerance of five points. The crossover is checked against an independent probe
with different sampling, which agrees at the middle horizons and disagrees at the
shortest one; the two curves are reported separately and never merged. The economic
conclusion is re-run over a grid of cost coefficients, which is why
Section~\ref{sec:p2disc} quotes a range, not a single optimum. Conditional coverage at the
longest official prefix ranges from \slotPConCovMin{} to \slotPConCovMax{} across the
\slotPConCovNEnv{} test environments, against a marginal \slotPConCovMarginal{}; one
environment falls below $0.85$ (\slotPConCovBelow{} of \slotPConCovNEnv{}). The marginal
guarantee does not promise the former, and with one to three test records per environment a
per-environment conformal calibration cannot be estimated, so we report the spread as a
limitation instead of calibrating four radii.

\noindent\textbf{Part III.} The paired comparison is the primary evidence: every policy sees the same durations on the
same seed, so seed noise cancels. The interval against the reference excludes zero in both
informative cases, so the ordering is not seed luck. Against \textsc{risk}, the variant that
spends reviews, the interval
\slotPStochPairSTwoRiskFusionDiff{}$\pm$\slotPStochPairSTwoRiskFusionCI{} and
\slotPStochPairSThreeRiskFusionDiff{}$\pm$\slotPStochPairSThreeRiskFusionCI{} covers zero,
so paying for observations is not a gain this experiment can detect. The fused-predictor
variant is worse than the plain one in S\slotCaseTwo{} by
\slotPStochPairSTwoSlackFusionDiff{} and indistinguishable in S\slotCaseThree, a warning
about stacking models, not a result about scheduling.

Two secondary checks bound the claim. The first is the mapping, and Table~\ref{tab:maps} puts
it on the record: across the five published mappings the reference moves by
\slotPMapRefSpanSTwo{} points in S\slotCaseTwo{} and \slotPMapRefSpanSThree{} in
S\slotCaseThree, while the frozen rule moves by \slotPMapOursSpanSTwo{} and
\slotPMapOursSpanSThree{}. No scheduler-only change in this paper moves the rate by that much,
so the mapping stays in the paper as the dominant factor, and the gain of the rule is quoted
inside a mapping and never across one. It is not uniform: the mean paired gain is
\slotPMapGainMeanSTwo{} and \slotPMapGainMeanSThree{} points, the weakest cell gains
\slotPMapGainMinSTwo{} and \slotPMapGainMinSThree{}, and \slotPMapSatSTwo{} cell in
S\slotCaseTwo{} is saturated, where the reference already reaches the rule.
